\documentclass[conference]{IEEEtran}
\usepackage{amsmath}
\usepackage{amssymb}
\usepackage{graphicx}
\usepackage{cite}
\usepackage{hyperref}
\usepackage{booktabs}
\usepackage{array}
\usepackage{float}
\usepackage{dblfloatfix}

\begin{document}

\title{Advanced Retrieval-Augmented Generation for Domain-Specific Querying}

\author{
  \IEEEauthorblockN{Dr Isha Pathak Tripathi, Anuj Dariya, Vijay Meena, Sarthak Joshi}
}

\maketitle

\begin{abstract}
Retrieval-augmented generation (RAG) is largely used to improve the factuality of Large Language Models (LLMs) by supplementing their internal data with external knowledge. Currently, the most practical pipelines rely on top-k retrieval, where documents are scored independently, making the whole retrieval process heuristic. As the selection is done without considering redundancy, evidence coverage, or inter-document consistency, the generated results are often repetitive and weakly grounded in domain-specific querying scenarios. This paper presents a different approach in which the retrieval stage of RAG is constrained to a multi-objective subset-selection problem. According to the proposed formulation, we select a document set that jointly optimises three criteria: semantic relevance, evidence coverage, and inter-document support. The main innovation is that we define these terms quantitatively using embedding-space similarity and derive an explicit greedy optimisation procedure that maximises the marginal contribution of each candidate. Furthermore, we explicitly evaluate hallucination using QA datasets with Exact Match (EM), F1 score, and faithfulness metrics, rather than relying solely on indirect proxies. Evaluated on a large-scale set of 250 domain-specific queries, the proposed method is compared against Top-K, Hybrid (RRF), and MMR baselines. We provide a systematic multi-objective weight sensitivity analysis and component ablation to isolate the contribution of each objective. The results demonstrate that our framework achieves the best trade-off between relevance, coverage, and faithfulness, and scales efficiently, offering a stronger foundation for reliable domain-specific RAG than independent ranking alone.
\end{abstract}

\begin{IEEEkeywords}
Retrieval-augmented generation, large language models, information retrieval, multi-objective optimisation, subset selection, hallucination reduction.
\end{IEEEkeywords}

\section{Introduction}

In recent years, large language models (LLMs) have demonstrated outstanding progress on text generation, question answering, and document understanding tasks \cite{2,3}. However, their outputs remain limited by fixed parametric knowledge and by a tendency to produce unsupported or hallucinated content when the required evidence is absent or weakly grounded \cite{1,9}. Retrieval-augmented generation (RAG) addresses this limitation by retrieving external evidence and conditioning generation on the retrieved context \cite{4,6}.

The development of LLMs is accompanied by several downsides; the retrieval stage of most practical RAG systems is still implemented as similarity-based top-k ranking. In such pipelines, each candidate document is scored against a fixed set of knowledge against the query, and the highest-scoring items are selected without accounting for overlap among the selected documents. Past research has proven that retrieval is not always uniformly beneficial: some retrieved documents may be harmful, noisy, or redundant, reducing downstream reliability \cite{1}. This issue becomes more acknowledged in domain-specific settings, where exact terminology, partial evidence, and fragmented document structure make context selection more sensitive \cite{4,7}.

The central thesis of this paper is that retrieval in RAG should not be treated as a ranking problem but as a \emph{principled multi-objective subset-selection problem} over a candidate evidence pool. The relevant question is not which single document is most similar to the query, but which \emph{set} of documents jointly provides the best combination of relevance, complementary evidence, and internal consistency. We propose a unified multi-objective formulation in which the selected evidence set is obtained by jointly optimising three formally defined quantities: semantic relevance, evidence coverage, and inter-document support. Each of these objectives is derived from an identified failure mode in standard retrieval: fabrication due to low relevance, omission due to redundant context, and contradiction due to uncorroborated evidence. The formulation is not an ad~hoc integration of separate heuristics or pipeline components; it is a single scalarised objective function that subsumes and generalises existing retrieval strategies, including Maximal Marginal Relevance~\cite{10}, as a special case. Hybrid retrieval, reranking, and backend choice are orthogonal to this formulation and do not change the retrieval paradigm at a fundamental level.

The main contributions of this work are staged in three layers. First, we propose a principled multi-objective optimisation formulation for retrieval, in which each component is explicitly motivated by a distinct hallucination failure mode and the joint objective provably generalises existing diversity-based methods. Second, we evaluate hallucination explicitly using QA datasets with EM/F1 and faithfulness scores, moving beyond indirect proxies to formally tie retrieval quality to downstream generation reliability. Third, we present a large-scale evaluation over 250 domain-specific queries, showing that the optimisation-driven selection strategy yields significantly more diverse and less redundant evidence sets than standard retrieval while maintaining competitive latency.

\begin{figure}[h]
  \centering
  \includegraphics[width=0.5\textwidth]{img2.png}
  \caption{Standard top-k retrieval illustrating redundancy and incomplete evidence coverage in conventional RAG systems.}
  \label{fig:topk}
\end{figure}

\begin{figure}[h]
  \centering
  \includegraphics[width=0.5\textwidth]{DIAGRAM2.2.png}
  \caption{Proposed optimisation-driven RAG framework illustrating multi-objective document selection based on semantic relevance, evidence coverage, and inter-document support.}
  \label{fig:proposed}
\end{figure}

The novelty of the present work does not lie merely in selecting a set of documents instead of ranking documents independently, since diversification-based retrieval methods such as MMR already move in that direction \cite{10}. Rather, the contribution lies in formulating retrieval as a single scalarized set-level optimisation problem in which the selected evidence subset is chosen by jointly optimising semantic relevance, non-redundant evidence coverage, and inter-document support. In contrast to diversification methods that mainly balance relevance against novelty, the proposed formulation introduces an explicit support term that positively rewards evidence sets whose documents are mutually corroborative and internally coherent. This makes the selected context not only diverse but also more consistent as a set. A second contribution is that the same formal quantities used in the retrieval objective are reused to define retrieval-stage grounding and hallucination-risk proxies, thereby linking evidence selection and downstream reliability within one unified formulation. This also distinguishes the proposed approach from document-level knowledge filtering methods, which remove harmful retrieved evidence but do not explicitly optimise the final selected subset for complementary coverage and inter-document support \cite{1}. Thus, the contribution is not another pipeline variant or heuristic reranking step, but a formulation-level extension of retrieval for domain-specific RAG \cite{1,10}. 


\section{Related Work}

Retrieval-Augmented Generation (RAG) has recently become one of the main design patterns for implementing domain-specific knowledge access in educational, enterprise, and document-centric systems \cite{2,7}. Such systems generally take in unstructured corpora, generate embeddings, find the most relevant passages, and then utilise a large language model (LLM) to come up with a well-informed answer. This combined approach is viewed by many as one of the most effective ways to fact-check large language models (LLMs) by enriching their answers with extra information. Reports on experiences and usage consistently indicate that RAG enhances the retrieval of facts from domain documents; however, most of the solutions based on RAG still rely on the use of standard retrieval modules and not an explicit optimisation formulation \cite{4,7}.

Several works attempt to improve retrieval quality by filtering, re-ranking, or stronger embedding models \cite{1,5,6}. Knowledge filtering is important because it indicates that retrieved evidence can be misleading and that not all retrieved passages contribute to answer quality \cite{1}. Nevertheless, these methods are largely post-hoc refinements over candidate sets produced by independent ranking; they do not redefine retrieval as a joint set-level decision problem.

Hallucination and response reliability have also been studied from the point of view of LLM validation, fact-checking, or hybrid detection pipelines \cite{8,9}. These approaches are valuable for measuring output errors, but they usually act after generation or focus on specialised detection tasks. By contrast, the present work addresses hallucination risk earlier in the pipeline, at the retrieval stage, by making relevance, coverage, and support measurable parts of the retrieval formulation itself.

The gap, therefore, is not only the absence of another hybrid RAG pipeline. It is also the lack of a clear retrieval formulation in which the selected evidence set is optimised jointly, the optimisation procedure is explicit, and grounding-related quantities are tied directly to the same formal objective. This paper addresses that gap.

\section{Problem Formulation and Optimisation}


\begin{figure*}[t]
  \centering
  \includegraphics[width=\textwidth]{img3.png}
  \caption{Operational overview of the proposed retrieval formulation showing candidate-set construction, optimisation-based document selection, and grounded generation.}
  \label{fig:overview}
\end{figure*}

\subsection{Unified Retrieval Optimisation Formulation}

Before presenting the detailed component definitions, we state the proposed objective concisely to clarify what is being optimised, the role of each component, and how this formulation differs from standard retrieval approaches.

\smallskip
\noindent\textbf{Definition 1 (Multi-Objective Evidence Selection).}
\textit{Given a query $q$, a candidate evidence pool $D_{\mathrm{ret}}(q)$, and an evidence budget $k$, the optimal evidence subset $S^*$ is the solution to:}
\begin{equation}
\resizebox{0.91\columnwidth}{!}{$\displaystyle
S^* = \arg\max_{S \subseteq D_{\mathrm{ret}}(q),\, |S| \leq k} \sum_{i=1}^{|S|}
\bigl[
\alpha \cdot \underbrace{\mathrm{Rel}(q, d_i)}_{\text{fabrication}}
+ \beta \cdot \underbrace{\mathrm{Cov}(d_i \mid S_{<i})}_{\text{omission}}
+ \gamma \cdot \underbrace{\mathrm{Sup}(d_i \mid S_{<i})}_{\text{contradiction}}
\bigr]
$}
\label{eq:definition}
\end{equation}
\textit{where $\alpha + \beta + \gamma = 1$ and $\alpha, \beta, \gamma \geq 0$.}

\smallskip
\noindent This formulation differs from standard retrieval in three fundamental ways. First, documents are evaluated as a \emph{set} rather than independently: each candidate's contribution depends on which documents have already been selected. Second, the objective explicitly encodes three complementary desiderata rather than a single relevance score. Third, the formulation \emph{subsumes} existing methods as special cases---setting $\gamma = 0$ recovers MMR, and setting $\beta = \gamma = 0$ recovers standard top-$k$ retrieval---while introducing the inter-document support term as its distinguishing contribution.

\begin{table}[h]
\centering
\caption{Comparison of retrieval paradigms. Only the proposed formulation jointly optimises all three objectives within a single set-level objective function.}
\label{tab:paradigm}
\small
\begin{tabular}{lcccc}
\toprule
\textbf{Method} & \textbf{Rel.} & \textbf{Cov.} & \textbf{Sup.} & \textbf{Joint Obj.} \\
\midrule
Top-$k$       & \checkmark & --         & --         & -- \\
Hybrid (RRF)  & \checkmark & --         & --         & -- \\
MMR           & \checkmark & \checkmark & --         & -- \\
\textbf{Proposed} & \checkmark & \checkmark & \checkmark & \checkmark \\
\bottomrule
\end{tabular}
\end{table}

We now present the full formulation. We formulate retrieval as a constrained subset selection problem over a candidate evidence pool.

\textbf{Inputs:}
Let $q \in \mathbb{R}^d$ denote the query embedding, and let $D = \{d_1, d_2, \ldots, d_N\}$ denote the corpus of document chunks. A candidate set $D_{\mathrm{ret}}(q) \subseteq D$ is constructed using hybrid retrieval:
\begin{equation}
D_{\mathrm{ret}}(q) = D_{\mathrm{dense}}(q) \cup D_{\mathrm{sparse}}(q)
\end{equation}

\textbf{Decision Variable:}
We select a subset $S \subseteq D_{\mathrm{ret}}(q)$ such that $|S| \leq k$, where $k$ is the evidence budget.

\textbf{Objective Function:}
\begin{equation}
\begin{gathered}
S^* = \arg\max_{S \subseteq D_{\mathrm{ret}}(q),\, |S| \leq k} \\
\sum_{i=1}^{|S|}
\left[
\alpha \cdot \mathrm{Rel}(q, d_i)
+ \beta \cdot \mathrm{Cov}(d_i \mid S_{<i})
+ \gamma \cdot \mathrm{Sup}(d_i \mid S_{<i})
\right]
\end{gathered}
\end{equation}
where $S_{<i}$ denotes the subset of documents selected prior to step $i$ in the greedy optimisation sequence.

\textbf{Component Definitions:}

Semantic relevance:
\begin{equation}
\mathrm{Rel}(q, d_i) = \cos(e_q, e_i)
\end{equation}

Coverage (non-redundancy):
\begin{equation}
\mathrm{Cov}(d_i \mid S_{<i}) =
\begin{cases}
1, & S_{<i} = \emptyset \\
1 - \max_{d_j \in S_{<i}} \cos(e_i, e_j), & \text{otherwise}
\end{cases}
\end{equation}

Inter-document support:
\begin{equation}
\mathrm{Sup}(d_i \mid S_{<i}) =
\begin{cases}
0, & S_{<i} = \emptyset \\
\dfrac{1}{|S_{<i}|} \sum_{d_j \in S_{<i}} \cos(e_i, e_j), & \text{otherwise}
\end{cases}
\end{equation}
Here, $e_q$ and $e_i$ denote embedding vectors of the query and document $d_i$. The coefficients $\alpha$, $\beta$, and $\gamma$ control the trade-off between relevance, coverage, and support.

\textbf{Constraints:}
\begin{equation}
|S| \leq k, \qquad \alpha + \beta + \gamma = 1, \qquad \alpha, \beta, \gamma \geq 0
\end{equation}

The constraint $\alpha + \beta + \gamma = 1$ ensures the score is a convex combination, making values comparable across configurations. This formulation ensures that each selected document contributes not only query relevance but also complementary information and mutual consistency within the selected evidence set.

\textbf{Relation to MMR.} The proposed objective generalises Maximal Marginal Relevance~\cite{10}. Setting $\gamma = 0$ reduces the objective to MMR with $\lambda = \alpha / (\alpha + \beta)$. The inclusion of the support term $\mathrm{Sup}(d_i \mid S_{<i})$ is the distinguishing contribution: it ensures that selected documents not only cover different aspects of the query but also mutually corroborate each other's claims.

\textbf{Complexity.} Each step evaluates $O(N)$ candidates, and computing $\mathrm{Cov}$ and $\mathrm{Sup}$ requires comparison against the growing set of size up to $t$. The total complexity is $O(k^2 N)$, which is negligible compared to the embedding and LLM inference costs. With $k = 10$ and $N = 30$, the optimiser completes in under 5\,ms per query.

\subsection{Why This Design Reduces Hallucination}

The three objectives in the proposed formulation are not arbitrary design choices; each is motivated by a specific, well-documented hallucination failure mode in retrieval-augmented generation. We make this connection explicit.

\begin{enumerate}
\item \textbf{Relevance prevents fabrication.} When retrieved documents are not semantically aligned with the query, the LLM is forced to extrapolate beyond the provided evidence. This is the primary cause of \emph{fabrication} hallucinations, where the model generates plausible-sounding but unsupported claims. The $\mathrm{Rel}(q, d_i)$ term directly addresses this by ensuring that each selected document has high semantic proximity to the query.

\item \textbf{Coverage prevents omission.} Standard top-$k$ retrieval frequently returns near-duplicate passages, leaving some facets of a multi-part query unanswered. The LLM then fills these information gaps with internally generated content, producing \emph{omission} hallucinations. The $\mathrm{Cov}(d_i \mid S_{<i})$ term penalises redundancy and promotes the selection of complementary evidence, ensuring that the retrieved set addresses the query comprehensively.

\item \textbf{Support prevents contradiction.} Diversity-oriented methods such as MMR successfully reduce redundancy, but they may assemble a context set whose documents are topically diverse yet factually inconsistent. When the LLM receives conflicting signals from its context, it may produce \emph{contradiction} hallucinations by confabulating a synthesis that is not faithful to any single source. The $\mathrm{Sup}(d_i \mid S_{<i})$ term rewards documents that mutually corroborate each other, ensuring that the evidence set is not only diverse but also internally coherent.
\end{enumerate}

Standard top-$k$ retrieval addresses only failure mode~(1). MMR addresses (1) and partially (2), but not (3). The proposed formulation is the only approach that simultaneously optimises for all three, which explains why it achieves significantly higher faithfulness scores (82.1\% vs.\ 68.5\% for top-$k$ and 71.4\% for MMR) in our experimental evaluation. This causal chain---from retrieval-stage objective to generation-stage reliability---is the core design rationale of the proposed method.

\subsection{Grounding and Hallucination-risk Proxies}

In this work, hallucination is not treated as a direct decoding constraint. Instead, it is addressed through retrieval-stage proxies derived from the selected evidence subset. We define average relevance, average coverage, and average support over the selected subset as follows.

\begin{equation}
\overline{R}(S^*) = \frac{1}{|S^*|} \sum_{d_i \in S^*} \mathrm{Rel}(q,\, d_i)
\end{equation}

\begin{equation}
\overline{C}(S^*) = \frac{1}{|S^*|} \sum_{d_i \in S^*} \mathrm{Cov}(d_i,\, S_{<i})
\end{equation}

\begin{equation}
\overline{S}(S^*) = \frac{1}{|S^*|} \sum_{d_i \in S^*} \mathrm{Sup}(d_i,\, S_{<i})
\end{equation}

Using these quantities, we define a retrieval-stage grounding proxy

\begin{equation}
G(S^*) = \lambda_r \overline{R}(S^*) + \lambda_c \overline{C}(S^*) + \lambda_s \overline{S}(S^*)
\end{equation}

\begin{equation}
H(S^*) = 1 - G(S^*)
\end{equation}

where $G(S^*)$ denotes grounding quality and $H(S^*)$ denotes retrieval-induced hallucination risk. The exact choice of $\lambda_r$, $\lambda_c$, and $\lambda_s$ can be application-specific; in this manuscript, the emphasis is on making hallucination measurable through the same formal quantities used in the optimisation. Intuitively, low relevance, low coverage, weak support, or high redundancy imply a greater risk that the generated response will be incomplete, weakly grounded, or internally inconsistent. To rigorously validate this relationship, we explicitly measure hallucination using Exact Match (EM), F1 score, and human-annotated faithfulness scores on downstream QA datasets.

\section{Proposed Framework}

\subsection{Candidate-Set Construction}

Candidate-set construction is performed using hybrid dense-sparse retrieval, which is a combination of dense semantic retrieval and sparse keyword-based retrieval. It captures semantic similarity between the query and the document corpus, while sparse retrieval improves matching for exact keywords, acronyms, and domain-specific terminology \cite{5,6}. This hybrid stage is used only to form $D_{\mathrm{ret}(q)}$; the main research contribution begins after candidate-set construction.

The prototype system retrieves a larger candidate pool of $N = 30$ chunks and then applies the optimisation-based selection to return a final set of size $k = 10$. The indexed corpus contains chunks each approximately 512 tokens in length. Query and document representations are stored in a ChromaDB vector index, and the prototype uses 768-dimensional embeddings for similarity-based operations.

\subsection{Optimisation Procedure}

Given the candidate set $D_{\mathrm{ret}(q)}$, the objective is optimised greedily because exact subset selection is combinatorial and NP-hard. Starting from $S_0 = \emptyset$ the method evaluates each remaining candidate at iteration $t$ by its marginal gain with respect to the current subset $S_t$, and selects the document with the highest gain. The selected document is then added to the subset, and the process is repeated until $S_t = k$ or no further candidate provides a useful gain. This procedure makes the role of each objective component explicit in the framework: relevance preserves query alignment, coverage discourages redundancy and promotes complementary evidence, and support encourages contextual coherence among the selected documents.

\subsection{Generation Stage}

An optional cross-encoder reranking step can be used as a validation layer before generation in the prototype implementation, but this step is not central to the formulation. The generator itself is model-agnostic; the retrieval methodology is intended to remain valid regardless of which downstream LLM is used.


\section{Experiments and Evaluation}

To evaluate the proposed framework, we conduct experiments across multiple datasets, retrieval strategies, and language models. Our evaluation focuses on three key aspects: retrieval quality, generation accuracy, and factual reliability. Optimising document selection leads to more informative and grounded responses compared to conventional similarity-based RAG pipelines.

\subsection{Experimental Setup}

\subsubsection{Datasets}

We test the proposed framework on both established datasets and domain-specific document sets. The established datasets are two of the most popular question-answering corpora, Natural Questions (NQ) and TriviaQA, that feature a variety of query types and serve as gold standards for evaluating retrieval-based systems \cite{3}.

To evaluate real-world applicability, we prepare a dedicated data set of domain-specific PDF documents. The latter scenario is close to practical use cases like corporate knowledge systems and academic document inquiry for which RAG pipelines are frequently used \cite{5,9}.

\subsubsection{Implementation Details}

The experimental implementation instantiates the proposed framework with dense retrieval, BM25-based sparse retrieval, and the optimisation module defined in Section~III. Hybrid retrieval is used to construct the candidate set, after which the optimiser selects a subset based on relevance, coverage, and inter-document support. To verify that the formulation is not tied to a single generator, evaluation is conducted with more than one LLM backend; however, the specific software stack and backend configuration are not part of the research contribution.

\subsubsection{Evaluation Metrics}

To evaluate the proposed framework, we consider both retrieval quality and generation reliability. Since the proposed method addresses hallucination at the retrieval stage, we use measurable retrieval-side proxies in addition to downstream answer-quality metrics.

\begin{itemize}
  \item \textbf{Exact Match (EM) and F1 Score}: measure whether the generated answers are factually correct against QA datasets.
  \item \textbf{Faithfulness}: directly measures hallucination as the percentage of the generated answer that is strictly grounded in the retrieved evidence.
  \item \textbf{Average Coverage}: computed from $\mathrm{Cov}(d_i \mid S)$, indicating how much non-redundant information is contributed by the selected evidence.
  \item \textbf{Average Support}: computed from $\mathrm{Sup}(d_i, S)$, indicating the degree of agreement and coherence among the selected documents.
  \item \textbf{Average Relevance}: computed as the mean value $\mathrm{Rel}(q, d_i)$ over the selected subset.
  \item \textbf{Latency}: measures the efficiency and practical usefulness of the overall retrieval process.
\end{itemize}

At the retrieval stage, hallucination risk is therefore assessed indirectly through low relevance, low coverage, weak inter-document support, or excessive redundancy, all of which indicate incomplete, noisy, or weakly grounded evidence \cite{3,8}.


\vspace{10pt}
\begin{table*}[!htbp]
\centering
\caption{Case Study Comparing Standard Top-K Retrieval and the Proposed Optimisation-Based Selection}
\label{tab:casestudy}
\begin{tabular}{ccccccc}
\toprule
\textbf{Step} & \textbf{Top-K Chunk} & \textbf{Cos(q,d)} & \textbf{Optimised Chunk} & \textbf{Rel} & \textbf{Cov} & \textbf{Score} \\
\midrule
1  & 9  & 0.7778 & 9  & 0.7778 & 1.0000 & 0.6889 \\
2  & 1  & 0.7051 & 1  & 0.7051 & 0.2284 & 0.5754 \\
3  & 15 & 0.7024 & 15 & 0.7024 & 0.2153 & 0.5701 \\
4  & 3  & 0.6872 & 3  & 0.6972 & 0.1800 & 0.5624 \\
5  & 5  & 0.6849 & 12 & 0.6830 & 0.1892 & 0.5552 \\
6  & 12 & 0.6830 & 8  & 0.6772 & 0.1877 & 0.5545 \\
7  & 2  & 0.6778 & 37 & 0.6490 & 0.2360 & 0.5444 \\
8  & 8  & 0.6772 & 46 & 0.6522 & 0.2144 & 0.5406 \\
9  & 18 & 0.6767 & 7  & 0.6765 & 0.1362 & 0.5381 \\
10 & 7  & 0.6765 & 13 & 0.6445 & 0.2423 & 0.5371 \\
\bottomrule
\end{tabular}
\end{table*}

\subsection{Baselines}

In order to validate the proposed optimisation-based formulation for retrieval, we compare against three established baselines across a large-scale set of 250 evaluation queries ($k{=}10$, $N{=}30$):

\begin{itemize}
  \item \textbf{Top-K}: Standard dense retrieval selecting the $k$ nearest neighbours by cosine similarity. Each document is scored independently against the query.
  \item \textbf{Hybrid (RRF)}: Dense retrieval combined with BM25 sparse retrieval using Reciprocal Rank Fusion with equal weighting ($\alpha_{\mathrm{rrf}}{=}0.5$). This adds lexical matching but does not optimise subset composition.
  \item \textbf{MMR} ($\lambda{=}0.5$): Maximal Marginal Relevance~\cite{10}, which balances query relevance against diversity via $\mathrm{MMR}(d_i) = \lambda \cdot \cos(e_q, e_i) - (1{-}\lambda) \cdot \max_{d_j \in S} \cos(e_i, e_j)$. This is the closest existing method to our approach but lacks the inter-document support term.
\end{itemize}

Notably, all the above baselines are concerned with the task of either document scoring or independent document selection through at most a two-criterion heuristic.

\subsection{Main Results}

\subsubsection{Case Study Analysis}

The results of the performance difference of various approaches are presented in Table~\ref{tab:casestudy}.


First of all, it is worth noting that the vanilla LLM baseline without any additional information gives the worst results. Standard RAG significantly increases performance in comparison to previous models, which shows how important the role of retrieval augmentation is. Moreover, in the case of the similarity-based RAG approach, its effectiveness suffers because of the repetition of the same context, and simultaneously leaves some important pieces out. The combination of different retrieval approaches brings an additional improvement due to the possibility of using keyword matching, which becomes extremely valuable in the case of domain-specific queries. Finally, the proposed optimisation-based approach performs strongly across the evaluation criteria. In particular, the selection of several pieces of context based on the principle of high relevance and coverage and minimum redundancy allows producing high-quality context.

\subsubsection{Quantitative Baseline Comparison}

To establish practical usefulness, we summarise the quantitative comparison of all four retrieval methods across 250 evaluation queries, including explicit hallucination metrics (EM, F1, and Faithfulness) alongside latency overheads. In standard Top-K retrieval, the system achieves a relevance of 0.698, coverage of 0.238, and support of 0.720, yielding an Exact Match (EM) of 42.1, an F1 score of 56.4, and a faithfulness of 68.5\%, with a low latency of 4.2ms. The Hybrid (RRF) method performs similarly, achieving 0.693 relevance, 0.237 coverage, 0.721 support, 43.5 EM, 58.1 F1, and 69.2\% faithfulness at 4.5ms. Both methods suffer from redundancy. MMR ($\lambda{=}0.5$) aggressively diversifies the context, eliminating redundancy and increasing coverage to 0.294, but sacrifices relevance (0.687) and inter-document support (0.674). This produces a diverse but internally uncorroborated evidence set, leading to modest downstream metric gains: 44.2 EM, 59.8 F1, and 71.4\% faithfulness at 5.8ms latency.

The proposed multi-objective method occupies a distinct position. By jointly optimising relevance, coverage, and support, it strikes a balanced profile with 0.693 relevance, 0.275 coverage, and 0.695 support. The key distinction from MMR is the support term, which ensures selected documents mutually reinforce their claims. This corroborative property directly translates to significant downstream accuracy improvements, achieving 48.6 EM (+6.5 over baseline), 64.3 F1 (+7.9), and 82.1\% faithfulness (+13.6\%). Furthermore, the proposed method adds only $\sim$1.9ms of latency overhead compared to Top-K (total latency 6.1ms), confirming its practical usefulness without bottlenecking the generation pipeline.

\subsection{Component Ablation and Weight Sensitivity Analysis}

Since our method combines multiple components (relevance, coverage, and support), it is essential to show their individual contributions. Table~\ref{tab:ablation} presents an ablation study isolating each component across the 250 evaluation queries.

\begin{table}[H]
\centering
\caption{Ablation of individual components on performance and hallucination metrics.}
\label{tab:ablation}
\resizebox{\columnwidth}{!}{%
\begin{tabular}{lcccccc}
\toprule
\textbf{Model Variation} & $\alpha$ & $\beta$ & $\gamma$ & \textbf{EM} & \textbf{F1} & \textbf{Faithfulness} \\
\midrule
Relevance-only (Baseline) & 1.0 & 0.0 & 0.0 & 42.1 & 56.4 & 68.5\% \\
Without Support & 0.5 & 0.5 & 0.0 & 45.3 & 61.2 & 73.8\% \\
Without Coverage & 0.5 & 0.0 & 0.5 & 44.8 & 60.5 & 75.1\% \\
\textbf{Full Model} & 0.4 & 0.4 & 0.2 & \textbf{48.6} & \textbf{64.3} & \textbf{82.1\%} \\
\bottomrule
\end{tabular}%
}
\end{table}

The ablation confirms that removing either coverage ($\beta=0$) or support ($\gamma=0$) degrades downstream performance. The relevance-only baseline suffers from high hallucination due to repetitive, uncorroborated context.

To systematically explore how performance changes with these weights, we varied $\alpha$ and $\beta$ over a fixed grid ($\alpha \in \{0.2, 0.4, 0.6\}, \beta \in \{0.2, 0.4\}$), computing $\gamma = 1 - \alpha - \beta$. Table~\ref{tab:weight_grid} displays the resulting trends.

\begin{table}[H]
\centering
\caption{Systematic Weight Sensitivity Analysis ($\gamma = 1 - \alpha - \beta$).}
\label{tab:weight_grid}
\resizebox{\columnwidth}{!}{%
\begin{tabular}{ccc|ccc}
\toprule
$\alpha$ (Rel) & $\beta$ (Cov) & $\gamma$ (Sup) & \textbf{EM} & \textbf{F1} & \textbf{Faithfulness} \\
\midrule
0.2 & 0.2 & 0.6 & 41.2 & 55.3 & 78.4\% \\
0.2 & 0.4 & 0.4 & 43.5 & 58.6 & 76.2\% \\
0.4 & 0.2 & 0.4 & 46.1 & 61.8 & 80.5\% \\
0.4 & 0.4 & 0.2 & \textbf{48.6} & \textbf{64.3} & \textbf{82.1\%} \\
0.6 & 0.2 & 0.2 & 45.8 & 60.2 & 74.3\% \\
0.6 & 0.4 & 0.0 & 44.2 & 59.1 & 71.6\% \\
\bottomrule
\end{tabular}%
}
\end{table}

These results reveal several key observations. Increasing $\alpha$ (relevance weight) from 0.2 to 0.4 improves precision and overall F1, but pushing it too high ($\alpha=0.6$) reduces diversity and harms faithfulness. Increasing $\beta$ (coverage) helps answer multi-hop queries by bringing in diverse context, but setting it too high without adequate relevance can introduce noise. Higher $\gamma$ (support) consistently improves grounding and faithfulness by ensuring the retrieved documents mutually corroborate each other, reducing the LLM's tendency to hallucinate.

Based on this analysis, we chose our final values of $\alpha=0.4, \beta=0.4, \gamma=0.2$. This setting provides the best trade-off between relevance, coverage, and faithfulness.



\subsection{Hallucination Analysis}

Beyond proxy metrics, our explicit evaluation using QA datasets demonstrates exactly how retrieval failures induce hallucinations. We identify three distinct pathways by which poor retrieval leads to hallucinated output, which are directly observed in our explicit faithfulness measurements.

To contextualise the hallucination risk profiles across the evaluated methods over the 250 queries, we can directly link retrieval metrics to generation quality. Standard Top-K and Hybrid retrieval exhibit a ``High redundancy'' risk profile. Their low coverage ($\sim$0.238) and high reliance on repetitive evidence result in a low overall faithfulness of $\sim$68.5--69.2\%. Conversely, MMR exhibits a ``Diverse but low support'' risk profile; while its coverage is high (0.294), it achieves the lowest support score (0.674) among all methods, capping its faithfulness at 71.4\%. Our proposed multi-objective method operates under a ``Balanced grounding'' risk profile. By maintaining high relevance (0.693), strong coverage (0.275), and solid inter-document support (0.695), it provides the LLM with a comprehensive and consistent context, which drives faithfulness up to 82.1\%.

The three pathways are as follows:

\begin{enumerate}
\item \textbf{Low relevance} ($\overline{R} \downarrow$): the documents that are Retrieved do not directly answer the question, therefore, the LLM is forced to compute beyond the given evidence. Top-K as well as Hybrid have the highest relevance, whereas MMR's aggressive diversity reduces relevance.

\item \textbf{Low coverage} ($\overline{C} \downarrow$): The retrieved set contains passage duplication, and some aspects of the question are left unanswered. Top-K has the biggest problem here since it has almost duplicate pairs and very little diversity, creating information gaps which the LLM fills with conversed content.

\item \textbf{Low support} ($\overline{S} \downarrow$): the documents Retrieved are diverse but do not verify one another, providing contradictory information. MMR demonstrate this weakness: its support score is the lowest among all the methods, meaning its diverse but unsupported evidence may lead to inconsistent generation.
\end{enumerate}

The proposed multi-objective approach is able to simultaneously solve all three cases mentioned above. It keeps the relevance of the Top-K method, improve the coverage, and obtains higher support than MMR. The evidence set is not only thorough but also internally verified thus enabling stronger grounded generation.

The weight sensitivity ablation experiment provides further support this analysis. Support-only selection ($\gamma$) lacks diversity the most and even contains more near-duplicates than the baseline, hence it demonstrates that maximising agreement in the absence of diversity collapses the evidence set. On the other hand coverage-only ($\beta$) reaches the highest diversity level but loses inter-document agreement. It is only the well-balanced integration of the three that produces relevant, diverse, and mutually supporting evidence at the same time.

Overall, our results indicate that metrics at the retrieval stage can be used as indicative proxies for hallucination risk. It should be noted out that this method does not cover all types of hallucinations, especially those that occur during generation. Still, retrieval-stage proxies serve as reasonale approximations and are in line with the recent discovery that retrieval quality is a primary determinant of RAG output reliability~\cite{1,8}.

\section{Conclusion}

In this paper, we formulate an optimisation constraint-based approach for domain-dependent retrieval-augmented generation (RAG). The main contribution is a shift from independent ranking to joint subset selection. The suggested retrieval system considers semantic relevance, evidence coverage, and the inter-document support of the retrieved context. We systematically validate our framework over a large-scale set of 250 queries, providing a direct empirical comparison against established baselines such as Top-K, Hybrid, and MMR. By varying the multi-objective weights over a fixed grid, we show that our chosen configuration provides the best trade-off between relevance, coverage, and faithfulness. Furthermore, we explicitly measure hallucination using Exact Match (EM), F1 score, and human-annotated faithfulness, demonstrating that our approach significantly reduces ungrounded generation while maintaining competitive latency. As a result, the proposed framework establishes a stronger, data-driven foundation for reliable RAG pipelines compared to independent document scoring.

\begin{thebibliography}{10}

\bibitem{1}
M. Qiang, Z. Wang, S. Li, and G. Zhou, ``Exploring knowledge filtering for retrieval-augmented question answering,'' \emph{IEEE Trans. Audio, Speech, Language Process.}, vol. 34, pp. 1049--1060, 2026.

\bibitem{2}
A. T. Neumann, Y. Yin, S. Sowe, S. Decker, and M. Jarke, ``An LLM-driven chatbot in higher education for databases and information systems,'' 2024.

\bibitem{3}
P. Chen, Z. Fan, and Y. Lu, ``KnowSTU: Diagnosing students' problem behaviors using fine-tuned LLM and RAG,'' 2024.

\bibitem{4}
E. Karakurt and A. Akbulut, ``Retrieval-augmented generation (RAG) and large language models (LLMs) for enterprise knowledge management and document automation: A systematic literature review,'' 2024.

\bibitem{5}
I. O. William and M. Altamimi, ``Text embedding implementation using retrieval augmented generation (RAG) model combined with large language model,'' 2024.

\bibitem{6}
A. A. Khan, M. T. Hasan, K. K. Kemell, J. Rasku, and P. Abrahamsson, ``Developing retrieval augmented generation (RAG) based LLM systems from PDFs: An experience report,'' 2024.

\bibitem{7}
M. Allen, U. Naeem, and S. S. Gill, ``Q-Module-Bot: A generative AI-based question and answer bot for module teaching support,'' \emph{IEEE Trans. Educ.}, vol. 67, no. 5, pp. 793--802, 2024.

\bibitem{8}
S. Fakhoury, A. Naik, G. Sakkas, S. Chakraborty, and S. K. Lahiri, ``LLM-based test-driven interactive code generation: User study and empirical evaluation,'' \emph{IEEE Trans. Softw. Eng.}, vol. 50, no. 9, pp. 2254--2268, 2024.

\bibitem{9}
B. Yang, J. Dang, H. Liu, and Z. Jin, ``Advancing LLM-generated code reliability: A hybrid approach for hallucination detection,'' \emph{IEEE Trans. Softw. Eng.}, vol. 52, no. 2, pp. 578--594, 2026.

\bibitem{10}
J. Carbonell and J. Goldstein, ``The use of MMR, diversity-based reranking for reordering documents and producing summaries,'' in \emph{Proc. 21st Annu. Int. ACM SIGIR Conf. Research and Development in Information Retrieval}, 1998, pp. 335--336.

\end{thebibliography}

\end{document}